\section*{第 \ref{ch:b3:glaciology} 章——冰川学}

\begin{solution}{exo:b3:glaciology:1}
$\tau_b = \rho gH\sin\alpha = 917\times9.81\times250\times\sin\ang{3} =
\qty{118}{kPa}$，约 \qty{1.2}{bar}，
与真实冰川下的应力量级相当。
\end{solution}

\begin{solution}{exo:b3:glaciology:2}
(a)~$\qty{2.9e15}{m^3}\times917 = \qty{2.66e18}{kg} = \qty{2.66e6}{Gt}$;
除以 \qty{360}{Gt/mm}，得 \qty{7390}{mm}，即 \qty{7.4}{m}，与所引数值一致。
(b)~$\qty{3.0e16}{m^3}\times917/360 = \qty{76}{m}$，对比 \qty{58}{m}。
格陵兰的冰大多搁在海平面以上的基岩上，因此其全部冰量
融化后都会使海洋增加水量。南极的冰大部分搁在海平面以下的
基岩上，这部分冰已经排开了海水，融化它只是补回被排开的水量；
漂浮的冰架则完全不增加水量。
\end{solution}

\begin{solution}{exo:b3:glaciology:3}
增量 $\num{6e6}\times1.2 = \qty{7.2e6}{m^3}$，损失 $\num{4e6}\times2.1 =
\qty{8.4e6}{m^3}$ 的水量，净值 \qty{-1.2e6}{m^3}，平均比物质平衡
为 $-1.2\times10^6/10^7 = \qty{-0.12}{m}$ 水当量。
\omterm{def:b3:glaciology:mass-balance}{积累区}面积比为 $6/10 = 0.6$。冰川正在失去质量，
即趋于萎缩。
\end{solution}

\begin{solution}{exo:b3:glaciology:4}
$1 - 820/917 = \num{0.106}$ 和 $1 - 830/917 = \num{0.095}$：约
十分之一的体积在孔隙封闭时仍为孔隙空间。
\end{solution}

\begin{solution}{exo:b3:glaciology:5}
$\tau_b = 917\times9.81\times300\times\sin\ang{3} = \qty{141.2}{kPa}$。由
板片定理，$u_s = \tfrac{2A}{4}\tau_b^3H = 0.5\times\num{1.43e-16}
\times(\num{1.412e5})^3\times300 = \qty{60}{m/\text{年}}$。厚度加倍时：$u_s
\propto H^{n+1} = H^4$，速度变为原来的十六倍。坡度减半时：$u_s \propto
(\sin\alpha)^3$，速度变为原来的八分之一。
\end{solution}

\begin{solution}{exo:b3:glaciology:6}
在 \qty{-25}{\celsius}（\qty{248.15}{K}）时，$A = \num{3.61e-13}\exp(-60000/
(8.314\times248.15)) = \qty{8.5e-26}{Pa^{-3}.s^{-1}}$，或者
\qty{2.7e-18}{Pa^{-3}.\text{年}^{-1}}。在 \qty{-10}{\celsius} 时，$\num{3.61e-13}
\exp(-27.42) = \qty{4.4e-25}{Pa^{-3}.s^{-1}} = \qty{1.4e-17}{Pa^{-3}.\text{年}^{-1}}$.
在 \qty{0}{\celsius}（暖支）时为 \qty{1.4e-16}{Pa^{-3}.\text{年}^{-1}}。
相同应力下，速度之比等于\omterm{def:b3:glaciology:glen}{速率因子}之比：
$\num{1.43e-16}/\num{2.67e-18} = 54$.
\end{solution}

\begin{solution}{exo:b3:glaciology:7}
$h = \sqrt{2\times10^5\times10^6/(917\times9.81)} = \qty{4.7}{km}$，接近
南极的 \qty{4.9}{km}：塑性模型仅凭冰盖的
宽度即可得到其厚度。
\end{solution}

\begin{solution}{exo:b3:glaciology:8}
$\tau = 200/5 = \num{40}\text{ 年}$。增暖停止后，冰川长度仍将持续调整
数十年（即若干个响应时间）。今天的
冰川尚未跟上今天的气候。
\end{solution}

\begin{solution}{exo:b3:glaciology:9}
(a)~$h = \sqrt{2\times2.03\times20\times\num{8.64e6}/(917\times\num{3.34e5})}
= \qty{1.51}{m}$。由于 $h^2 \propto t$，厚度增至 \qty{3}{m} 需要 $100\times
(3/1.51)^2 = 393$ 天，超过一个冬季。(b)~远厚于
\qty{3}{m} 的冰不可能在一个冬季内靠冻结增长而成：这种冰是
经历数个冬季存留下来的，或是由风和洋流在冰脊中堆积而成的冰。
\end{solution}

\begin{solution}{exo:b3:glaciology:10}
(a)~$H = (1025/917)\times800 = \qty{894}{m}$。(b)~当 \qty{900}{m} 时，
\omterm{def:b3:glaciology:grounding-line}{接地线}处厚度为 \qty{1006}{m}，为原来的 $9/8$；
通量增为 $(9/8)^5 = 1.8$ 倍。(c)~\omterm{def:b3:glaciology:grounding-line}{接地线}沿向内加深的床面后退时，
\omterm{def:b3:glaciology:grounding-line}{接地线}处冰层更厚，排泄更多，使后方冰层变薄，
使\omterm{def:b3:glaciology:grounding-line}{接地线}进一步后退：这是一种正反馈。冰架
支撑接地冰，可减小\omterm{def:b3:glaciology:grounding-line}{接地线}处的通量，并
即使在向下倾斜的床面上，也能将其稳定在原处。
\end{solution}

\begin{solution}{exo:b3:glaciology:11}
设 $h(r) = c\sqrt{R - r}$，$c = \sqrt{2\tau_0/(\rho g)}$，$V =
\int_0^R2\pi rc\sqrt{R - r}\,\mathrm dr$。令 $s = R - r$，
\[
\int_0^R(R - s)\sqrt s\,\mathrm ds = \tfrac23R^{5/2} - \tfrac25R^{5/2} =
\tfrac4{15}R^{5/2},
\]
于是 $V = \frac{8\pi}{15}cR^{5/2}$。$R$ 加倍时，
$V$ 与\omterm{def:b3:glaciology:sle}{海平面当量}均增为 $2^{5/2} = 5.7$ 倍。
\end{solution}

\begin{solution}{exo:b3:glaciology:12}
(a)~三月：$14.13/16.34 - 1 = \qty{-13.5}{\percent}$；九月：$4.81/7.05
- 1 = \qty{-31.8}{\percent}$。(b)~三月时，海冰边缘主要由
环绕北冰洋的海岸所控制，每年冬季的再冻结
使海盆重新充满；近年较薄的冰每年夏季
融化得更彻底。冰–反照率反馈作用于夏季：开阔水面
吸收约 \qty{94}{\percent} 的太阳辐射，海冰反射
\qtyrange{50}{70}{\percent}，因此夏季每损失一部分海冰，都会使海水升温，
进而融化更多海冰。
\end{solution}

\begin{solution}{pb:b3:glaciology:1}
\textbf{1.} 床面以上的冰柱，单位面积重量为 $\rho gH$，受
沿坡向下的分量 $\rho gH\sin\alpha$ 推动，由底部切应力平衡；
小坡度时 $\sin\alpha \approx \alpha$（
即以表面坡度代入的板片定理）。

\textbf{2.} $917\times9.81\times580\times0.0190 = \qty{99.1}{kPa}$;
$917\times9.81\times470\times0.0240 = \qty{101.5}{kPa}$; $917\times9.81
\times335\times0.0330 = \qty{99.4}{kPa}$.

\textbf{3.} 应力在 \qty{2}{\percent} 以内保持不变，而
厚度几乎减半：冰盖接近塑性，屈服应力约为 $\tau_0 \approx
\qty{100}{kPa}$。

\textbf{4.} $2\times10^5/(917\times9.81) = \qty{22.23}{m}$。

\textbf{5.} $h = \sqrt{22.23\,(R - r)}$：\qty{577}{m}、\qty{472}{m} 和
\qty{333}{m}，分别位于距边缘 15、10 和 \qty{5}{km} 处，
与测深值相差在数米以内。

\textbf{6.} $\sqrt{22.23\times20000} = \qty{667}{m}$。

\textbf{7.} 同 \cref{exo:b3:glaciology:11}：$V = \int_0^R2\pi r
\sqrt{22.23(R - r)}\,\mathrm dr = \frac{8\pi}{15}\sqrt{2\tau_0/(\rho
g)}\,R^{5/2}$。

\textbf{8.} $\frac{8\pi}{15}\times4.715\times(\num{2e4})^{5/2} =
\qty{4.47e11}{m^3} = \qty{447}{km^3}$。

\textbf{9.} $\pi\times20^2 = \qty{1257}{km^2}$；平均厚度
$447/1257 = \qty{0.356}{km}$，即 \qty{356}{m}，略大于
中心厚度的一半。

\textbf{10.} $2^{5/2} = 5.7$ 倍。

\textbf{11.} $A = \num{3.61e-13}\exp(-60000/(8.314\times263.15)) =
\qty{4.45e-25}{Pa^{-3}.s^{-1}}$，乘以 \qty{3.156e7}{s}：
\qty{1.40e-17}{Pa^{-3}.\text{年}^{-1}}。

\textbf{12.} 在 \qty{10}{km} 处，$H = \qty{471.5}{m}$，坡度 $471.5/20000 =
\num{0.0236}$，$\tau_b = \qty{100}{kPa}$；$u_s = \tfrac{2A}{4}\tau_b^3H =
0.5\times\num{1.40e-17}\times10^{15}\times471.5 = \qty{3.3}{m/\text{年}}$。

\textbf{13.} $q = \tfrac{2A}{5}\tau_b^3H^2 = 0.4\times\num{1.40e-17}
\times10^{15}\times471.5^2 = \qty{1250}{m^2/\text{年}}$；平均速度 $q/H =
\qty{2.6}{m/\text{年}}$，为表面速度的五分之四。

\textbf{14.} $2\pi rq = \pi r^2\bar b$，故 $\bar b = 2q/r = 2\times1250/10^4
= \qty{0.25}{m}$ 的冰（每年），分布于内侧 \qty{10}{km} 范围内。

\textbf{15.} 当 $h = 450$ 时，$R - r = 450^2/22.23 = \qty{9.11}{km}$，位于
$r = \qty{10.9}{km}$ 处；\omterm{def:b3:glaciology:mass-balance}{积累区}为其内的圆盘，面积占冰盖的 $(10.9/
20)^2 = 0.30$。

\textbf{16.} $A$ 增大 \num{10.2} 倍，相同
应力下的速度亦随之增大。冰的流出速度快于平衡所能补给的量：冰盖
将变薄并向外扩展，直至应力降低，重新与平衡相匹配；
较软的塑性冰盖屈服应力较小，
剖面也更平缓。

\textbf{17.} $\num{4.47e11}\times917 = \qty{4.10e14}{kg} = \qty{410}{Gt}$。

\textbf{18.} $410/360 = \qty{1.14}{mm}$。

\textbf{19.} $356/0.5 = 711$ 年。

\textbf{20.} 随着表面降低，冰盖处于更暖的空气中，融化加快，
即高程–质量平衡反馈；随着冰层变薄，应力和
流动也随之改变。变薄过程加速。

\textbf{21.} $410/243 = \num{1.7}\text{ 年}$：整个冰盖的质量
不足格陵兰岛在 2010--2019 年间两年的损失量。

\textbf{22.} 阿克尔冰盖的\omterm{def:b3:glaciology:sle}{海平面当量}为
\textbf{\qty[text-series-to-math=true]{1.1}{mm}}（\qty{1.14}{mm}）。
\end{solution}
